\hmsubsection{Technology Stack}{FS}{}\label{sec:technology-stack}
The technology stack was selected according to two criteria. First, the
tools had to run on the target hardware without custom drivers. Second,
each team member had to be able to set up the project on their own
laptop without external assistance.
The complete list is shown in Table \ref{tab:tech-stack}.
\begin{table}[h]
\centering
\small
\begin{tabular}{p{3.5cm} p{4cm} p{6.5cm}}
\toprule
\textbf{Component} & \textbf{Tool / library} & \textbf{Role} \\
\midrule
Host operating system & Raspberry Pi OS (64-bit) & Linux base for the Pi 5 node \\
Host language & Python 3.11 & Vision pipeline, inverse kinematics, cycle
orchestration \\
Image processing & OpenCV 4.x & HSV segmentation, Hough transform, contour analysis
\\
Camera driver & DepthAI v3 (Luxonis) & Interface to the OAK-D S2 \\
Numerics & NumPy, SciPy & Array operations and the Kalman filter used in the
tracker \\
Classifier (evaluated) & scikit-learn & SVM colour verification, later disabled \\
Firmware language & C++ (Arduino flavour) & Motor controller firmware on the
OpenRB-150 \\
Firmware libraries & Dynamixel2Arduino & Dynamixel Protocol 2.0 on the TTL bus \\
IDE (firmware) & Arduino IDE 2.x & Build and flash for the OpenRB-150 \\
Version control & Git, GitHub & Source repository with daily commits \\
Documentation & \LaTeX{} (Overleaf) & Single-source build of this report \\
Project tracking & Trello & Sprint backlog, in-progress, testing, done lanes \\
Communication & Slack, Microsoft Teams & Internal chat and supervisor meetings \\
\bottomrule
\end{tabular}
\caption{Technology stack used in the project.}
\label{tab:tech-stack}
\end{table}
Two choices require brief justification.
The host code was implemented in Python rather than C++. The vision
pipeline consists of a limited number of OpenCV and DepthAI calls. A
Python implementation requires approximately half as many lines of code
as a comparable C++ implementation, while the runtime is similar because
the computationally intensive operations are executed inside OpenCV,
which is written in C. Transitioning to C++ was therefore not expected to
improve runtime performance and would have reduced code readability.
The Arduino IDE was used for the firmware instead of PlatformIO. The
OpenRB-150 is officially supported through Arduino libraries from
ROBOTIS. A more advanced build system was unnecessary for the limited
firmware scope in this project. If the firmware expands, migration to
PlatformIO can be considered as a future improvement.
